\documentclass[9pt,a4paper]{article}

\usepackage[utf8]{inputenc}
\usepackage[T1]{fontenc}
\usepackage{lmodern}
\usepackage[french]{babel}
\usepackage[margin=1.2cm]{geometry}
\usepackage{microtype}
\usepackage{multicol}
\usepackage{enumitem}
\usepackage[table]{xcolor}
\usepackage[most]{tcolorbox}
\usepackage{amsmath}
\usepackage{amssymb}
\usepackage{eurosym}

\definecolor{navy}{HTML}{1F3864}
\definecolor{lightblue}{HTML}{D6E0F0}
\definecolor{warm}{HTML}{C55A11}
\definecolor{vert}{HTML}{1E6B3A}

\newtcolorbox{alerte}[1][]{colback=orange!8,colframe=orange!70!black,boxrule=0.6pt,
  arc=2pt,left=6pt,right=6pt,top=4pt,bottom=4pt,#1}
\newtcolorbox{reflexe}[1][]{colback=lightblue!50,colframe=navy,boxrule=0.8pt,
  arc=2pt,left=6pt,right=6pt,top=4pt,bottom=4pt,#1}
\newtcolorbox{vertb}[1][]{colback=vert!7,colframe=vert!70!black,boxrule=0.6pt,
  arc=2pt,left=6pt,right=6pt,top=4pt,bottom=4pt,#1}

\newcommand{\slide}[2]{%
  \medskip\noindent{\color{navy}\bfseries #1.}\ \textcolor{navy}{\footnotesize #2}\par\nobreak\smallskip}
\newcommand{\dire}[1]{\noindent\hangindent=1em\hangafter=1 #1\par}
\newcommand{\n}[1]{\textbf{#1}}
% Une definition : le terme, puis la phrase a dire telle quelle.
\newcommand{\dfn}[2]{\noindent\hangindent=1.1em\hangafter=1
  {\color{warm}\bfseries #1} \textemdash{} #2\par\smallskip}

\setlength{\parindent}{0pt}
\setlength{\columnsep}{14pt}
\pagestyle{empty}

\begin{document}

{\color{navy}\Large\bfseries Aide-mémoire, point tuteur du 7 août 2026}\\[2pt]
{\footnotesize Kélian Kaddouri. À garder sous les yeux. \textbf{Partie 1 : les définitions.}
\textbf{Partie 2 : slide par slide.} Les chiffres en gras sont ceux qu'il ne faut pas se tromper.}

\vspace{4pt}
\begin{alerte}
\footnotesize
\textbf{Ce qui s'est passé la dernière fois, et qui ne doit pas se reproduire.} Je n'ai pas su
définir \textbf{placebo} quand on me l'a demandé, alors que le mot était sur ma propre slide.
Un terme que je place et que je ne peux pas définir en une phrase annule le crédit de la
slide entière, et de celles d'avant. \textbf{Règle : je ne prononce aucun mot de la partie 1
sans pouvoir en dire la phrase à voix haute.} S'il en reste un que je ne tiens pas, je le
retire de mon vocabulaire du jour et je dis la chose en français.
\end{alerte}

\vspace{3pt}
{\color{navy}\large\bfseries Partie 1 \textemdash{} Les définitions}\\
{\footnotesize Chaque entrée se lit en une phrase. La phrase est \emph{suffisante} : si on
m'en demande plus, c'est écrit en italique.}

\vspace{2pt}
\begin{multicols}{2}
\scriptsize

{\color{navy}\bfseries A. Les tests, et d'abord celui que j'ai raté}

\dfn{Test placebo}{je casse volontairement la structure que je prétends détecter, je relance
\emph{le même} calcul, et je vérifie que le signal disparaît. S'il survit, ce n'est pas la
structure qui le produisait. \emph{Ici : je permute l'ordre du temps à l'intérieur de chaque
entité, ce qui détruit « qui vient avant qui » sans rien changer d'autre, ni les volumes ni
les co-occurrences.}}

\dfn{Loi nulle (ou loi de permutation)}{la distribution de ma statistique dans le monde où il
n'y a rien à trouver. Je l'obtiens en rejouant le placebo un grand nombre de fois. C'est mon
étalon : sans elle, un chiffre observé ne veut rien dire.}

\dfn{$z$ (score standardisé)}{de combien d'écarts-types mon observation s'écarte de la
moyenne de la loi nulle. $z=(\text{observé}-\text{moyenne du nul})/\text{écart-type du nul}$.
\emph{Repères : $|z|<2$ = indistinguable du hasard ; $z>3$ = fort.}}

\dfn{Mon $z=-0{,}33$}{asymétrie observée $119$, loi nulle $128\pm27$, donc
$(119-128)/27=-0{,}33$. \n{Ce n'est pas un échec, c'est le résultat} : sur une donnée sans
ordre temporel, la direction \emph{ne peut pas} se voir, et le placebo le confirme.}

\dfn{Mon $z=+3{,}9$ (post-mortems) et $+5{,}1$ (corpus étendu)}{là, le placebo est franchi :
le corpus documentaire, lui, porte l'ordre des événements.}

\dfn{$p$-value}{probabilité d'observer un résultat au moins aussi extrême \emph{si} il n'y
avait rien. Petite $=$ incompatible avec le hasard. \emph{Elle ne dit pas que l'effet est
grand, seulement qu'il n'est pas nul.}}

\dfn{Bootstrap}{je rééchantillonne mes données avec remise, je refais l'estimation à chaque
fois, et la dispersion des résultats me donne l'incertitude. \emph{Aucune hypothèse de loi.}}

\dfn{Bootstrap paramétrique}{même idée, mais je tire les échantillons \emph{depuis mon modèle
ajusté} au lieu des données. \emph{Nécessaire quand les valeurs critiques tabulées d'un test
ne valent plus, parce que les paramètres ont été estimés.}}

\dfn{Anderson-Darling, Kolmogorov-Smirnov}{deux tests d'adéquation : la loi que j'ai choisie
colle-t-elle aux données. AD pèse davantage la \emph{queue}, KS le milieu. \emph{Chez moi,
$p=0{,}96$ et $p=0{,}89$ : non rejetés.}}

\dfn{Jackknife}{je retire une observation, je recalcule, je recommence pour chacune. Ça
mesure l'influence de chaque point. \emph{Sert à vérifier qu'aucun résultat ne tient à un
seul incident.}}

\dfn{Knockout}{je retire un ingrédient du modèle et je regarde ce qui tombe. \emph{Ici : je
supprime l'asymétrie de $W$ et le signal passe sous le bruit du placebo, donc c'est bien
l'asymétrie qui porte le résultat.}}

\vspace{3pt}
{\color{navy}\bfseries B. La sévérité : la théorie des valeurs extrêmes}

\dfn{EVT}{la statistique des \emph{maxima}, là où le théorème central limite est celle des
sommes. Elle autorise à extrapoler au-delà du plus gros sinistre observé.}

\dfn{POT (\emph{peaks over threshold})}{plutôt que de garder un maximum par bloc, je garde
\emph{tout} ce qui dépasse un seuil $u$ et je modélise les excès. \emph{Moins de gaspillage
d'information, ce qui compte quand la donnée est rare.}}

\dfn{GPD (Pareto généralisée)}{la loi vers laquelle converge la distribution des excès quand
le seuil monte. \emph{C'est le théorème de Pickands-Balkema-de Haan qui le garantit.}}

\dfn{$\xi$ (indice de queue)}{le paramètre de forme de la GPD, celui qui commande la
lourdeur. \n{$\xi>0$ = queue lourde} ; la variance n'existe plus si $\xi\ge0{,}5$,
\n{l'espérance plus si $\xi\ge1$}. \emph{Chez moi $\xi=0{,}595$ sur OpRisk : variance
infinie, espérance finie.}}

\dfn{$\sigma$ (échelle GPD)}{le paramètre d'échelle des excès, en millions d'euros. \n{Ne pas
le confondre avec l'écart-type d'un test.} \emph{Chez moi $57{,}97$~M\euro.}}

\dfn{$u$ (seuil) et $p_u$ (taux de dépassement)}{$u=20{,}03$~M\euro{} et $p_u=\mathbb{P}(X>u)$.
\n{$p_u$ n'est pas libre : une fois le seuil et l'échantillon fixés, il se compte.}}

\dfn{Estimateur de Hill}{un estimateur non paramétrique de $\xi$, fondé sur les $k$ plus
grandes observations. \n{Il suppose une queue de Pareto pure et il est biaisé en échantillon
fini.}}

\dfn{MLE (maximum de vraisemblance)}{je choisis les paramètres qui rendent mes données les
plus probables. \emph{C'est l'estimateur que je retiens, sur les excès, parce qu'il n'exige
pas la Pareto pure.}}

\dfn{Information de Fisher}{la courbure de la vraisemblance à son sommet. Plus elle est
grande, plus l'estimation est précise. \emph{Elle me donne l'écart-type asymptotique, ici
$(1+\xi)/\sqrt{N_u}=0{,}167$.}}

\dfn{VaR $99{,}5\,\%$}{le montant dépassé une année sur deux cents. \n{C'est la définition
même du SCR sous Solvabilité~II.}}

\dfn{TVaR (\emph{expected shortfall})}{la perte \emph{moyenne} sachant qu'on a dépassé la
VaR. \emph{Elle regarde au-delà du quantile, la VaR s'arrête dessus.}}

\vspace{3pt}
{\color{navy}\bfseries C. La fréquence}

\dfn{Poisson}{la loi du nombre d'événements quand ils arrivent indépendamment à taux
constant. \emph{Sa propriété : variance $=$ moyenne.}}

\dfn{Surdispersion}{quand la variance observée dépasse la moyenne. \n{C'est le signe qu'un
facteur commun ou une hétérogénéité manque au modèle.}}

\dfn{Indice de dispersion $\varphi$}{le rapport variance sur moyenne. \n{Il n'est pas
invariant d'échelle} : il vaut $1+\lambda/r$, donc il croît avec le niveau de fréquence
auquel on le lit. \emph{D'où $1{,}19$ sur les cellules firme-année, $2{,}24$ sur la série
annuelle, $9{,}20$ au moteur qui est un choix posé.}}

\dfn{Binomiale négative}{un mélange de Poisson par une loi Gamma. \emph{C'est la façon
standard d'absorber la surdispersion sans changer la moyenne.}}

\dfn{Mélange Gamma, et pourquoi pas Erlang}{la forme ajustée vaut $0{,}63$, or une Erlang
exige une forme entière $\ge1$. \n{Le mélange Gamma est le bon objet, l'Erlang n'en est qu'un
cas particulier.}}

\dfn{Charge systémique $a$}{l'intensité du choc commun à tous les piliers, dans
$m=\exp(aY-a^2/2)$ avec $Y$ gaussien centré réduit. \emph{Le $-a^2/2$ est là pour que
$\mathbb{E}[m]=1$ : le choc redistribue, il n'ajoute pas de fréquence moyenne.}}

\dfn{Élasticité}{de combien de pour cent bouge une grandeur quand une autre bouge d'un pour
cent. \emph{Chez moi, sévérité contre taille : $0{,}087$, donc presque plate. C'est ce qui
limite la transposition vers les petites entités.}}

\dfn{Troncature}{les pertes sous le seuil de collecte n'existent pas dans la base.
\n{L'ignorer biaise l'estimation}, et le corriger est le rôle de l'ajustement tronqué.}

\vspace{3pt}
{\color{navy}\bfseries D. La cascade}

\dfn{$W$ (matrice de contagion dirigée)}{$W_{jk}$ est la force du lien \n{du pilier $j$ vers
le pilier $k$}. \n{La ligne émet, la colonne reçoit.}}

\dfn{Normalisation de Leontief}{je divise $\mathrm{TRANS}$ par l'\n{émission maximale} (2,60)
et je multiplie par $g$. \emph{Ça garantit que le pilier le plus prolifique engendre au plus
$g$ descendants directs, donc que le système est stable par construction et non par
hypothèse.}}

\dfn{$g$ (gain de propagation)}{la part de contagion du pilier le plus prolifique, dans
$(0,1]$. \n{Posé, non calibré}, et la thèse n'en dépend pas.}

\dfn{Rayon spectral $\rho(W)$}{la plus grande valeur propre en module. \n{Il gouverne la
stabilité} : si $\rho(W)<1$, la série $I+W+W^2+\cdots$ converge et $(I-W)^{-1}$ existe.
\emph{Chez moi $\n{0{,}506}$ à $g=0{,}9$, donc largement sous-critique.}}

\dfn{Sous-criticité}{la propriété $\rho<1$. \n{Elle veut dire que la cascade s'éteint.}}

\dfn{Processus de branchement}{un incident sur un pilier engendre, à la période suivante, un
nombre aléatoire d'incidents sur les autres. \emph{Toute la théorie de l'extinction des
épidémies s'applique.}}

\dfn{$R_0$}{le nombre moyen d'incidents engendrés par un incident, en cascade. \n{La cascade
s'éteint si et seulement si $R_0<1$.} \emph{Chez moi $\n{0{,}054}$ : c'est le $R_0$ des
épidémiologistes.}}

\dfn{$\rho(W)$ contre $R_0$, à ne pas confondre}{$\rho(W)$ gouverne le système latent
\emph{simultané}, $R_0$ la cascade d'\emph{incidents}. \emph{Les deux sont stables sur tout
le domaine admissible, mais pas du même ordre : la forme simultanée approche son seuil
critique plus de quatre fois plus vite.}}

\dfn{Identification partielle}{la donnée ne détermine pas le paramètre, mais elle le
\emph{borne}. \n{Je publie l'ensemble des valeurs compatibles, pas un point.} \emph{Ici : la
co-occurrence identifie la partie symétrique de $W$, jamais la direction.}}

\dfn{$W=S+A$}{toute matrice se coupe en une partie symétrique $S$ (qui est lié à qui) et une
partie antisymétrique $A$ (dans quel sens). \n{La donnée voit $S$, elle est aveugle à $A$.}}

\dfn{$t$ (degré d'ignorance sur la direction)}{la largeur que je m'autorise autour de $S$ :
$|A_{jk}|\le t\,S_{jk}$. \emph{$t=1$ = ignorance totale, et c'est l'état actuel.}}

\dfn{Sommets ($2^{10}=1024$)}{le maximum d'une fonction convexe sur un pavé est atteint sur un
coin. \emph{Je les énumère tous : les bornes sont exactes, pas approchées.}}

\vspace{3pt}
{\color{navy}\bfseries E. Le capital et l'allocation}

\dfn{SCR}{le capital réglementaire de Solvabilité~II, défini comme la VaR à $99{,}5\,\%$ de
la perte à un an.}

\dfn{ORSA}{l'évaluation interne des risques, pilier 2. \n{Il n'existe aucun module DORA en
Formule Standard} : ce que je calcule est un \n{besoin de capital ORSA}, jamais « le SCR
DORA » d'une entité.}

\dfn{Formule Standard, charge opérationnelle}{le forfait réglementaire, $3\,\%$ des
provisions, plafonné à $30\,\%$ du BSCR. \n{Il ne bouge pas d'un euro selon le profil de
conformité, et c'est tout mon argument.}}

\dfn{Marge de risque}{le coût de portage du capital, $6\,\%$ par an sous Solvabilité~II.
\emph{C'est ce qui transforme un $\Delta$SCR en euros annuels.}}

\dfn{$\Delta_{\text{DORA}}$}{l'écart de capital entre l'état non conforme et l'état conforme.
\n{C'est ma grandeur centrale, un écart, pas un niveau.}}

\dfn{Socle}{le SCR à contagion nulle ($W=0$). \n{Le plancher défendable} : tout ce qui est
au-dessus est de la contagion, donc à borner.}

\dfn{Super-additivité}{la somme des effets isolés est plus petite que l'effet conjoint.
\emph{Chez moi l'interaction vaut $5\,001$~M\euro, soit $35\,\%$ : les canaux ne
s'additionnent pas.}}

\dfn{Valeur de Shapley}{la seule façon d'attribuer un surcoût conjoint à des contributeurs
qui respecte l'efficacité (la somme fait le total), la symétrie et la linéarité. \emph{Vue
« source » : qui a causé le surcoût.}}

\dfn{Allocation d'Euler}{la contribution marginale de chaque brique au capital, valide parce
que la VaR est homogène de degré~1. \emph{Vue « réceptacle » : où le capital est consommé.}}

\dfn{Copule}{la fonction qui relie des lois marginales pour en faire une loi jointe.
\n{Elle décrit la dépendance sans toucher aux marges.}}

\dfn{Dépendance de queue (Gumbel contre gaussienne)}{la probabilité que deux risques soient
extrêmes \emph{ensemble}. \n{La gaussienne l'annule à la limite}, la Gumbel la conserve.
\emph{C'est pourquoi une copule gaussienne raterait un choc comme MOVEit.}}

\dfn{VaR prédictive}{j'intègre l'incertitude d'estimation \emph{avant} de prendre le
quantile. \emph{Le point honnête.}}

\dfn{VaR robuste}{je prends la borne haute sur l'ambiguïté du modèle. \emph{Le point
prudent. Le choix entre les deux est prudentiel, pas mathématique.}}

\dfn{Plug-in}{je prends le quantile de la loi estimée en un point, sans son incertitude.
\n{C'est ce qu'il ne faut pas reporter seul.}}

\end{multicols}

\vspace{2pt}
\begin{vertb}
\scriptsize
\textbf{Les six chiffres que je dois pouvoir sortir sans notes.}
$\rho(W)=\n{0{,}506}$ \textbullet{} $R_0=\n{0{,}054}$ \textbullet{} $\xi=\n{0{,}595}$ sur
$\n{91}$ excès \textbullet{} socle $\n{5\,275}$~M\euro{} \textbullet{} bande
$\n{[6\,858;8\,697]}$ \textbullet{} entité $\n{169}$~M\euro{} contre $\n{450}$ de forfait.
\end{vertb}

\newpage

{\color{navy}\large\bfseries Partie 2 \textemdash{} Slide par slide}

\vspace{2pt}
\begin{reflexe}
\textbf{La phrase à placer deux fois, minimum} (slide 8, puis en clôture)~:\\
\textbf{\og Le niveau bouge dans la bande, l'ordre des piliers ne bouge pas. C'est ça, mon
résultat.\fg{}}\\[3pt]
{\footnotesize Sans elle, la liste des incertitudes passe pour du flou. Avec elle, elle passe
pour de la rigueur.}
\end{reflexe}

\vspace{2pt}
\begin{multicols}{2}
\footnotesize

\slide{1. Titre}{}
Point d'avancement. Premier point où je donne un capital chiffré.

\slide{2. Depuis le 31 juillet}{annoncer le plan}
\dire{$\bullet$ Le placebo et la clarification $z$ / $\sigma$ sont partis par écrit, je n'y
reviens pas. \n{Mais je dois pouvoir les redéfinir si on me relance} (partie 1, section A).}
\dire{$\bullet$ Trois temps~: \n{(A)} le capital et son échelle, \n{(B)} tient-il face au
réel, \n{(C)} les vérifications demandées.}

\slide{3. Le capital et la contagion}{la table des bornes}
\dire{$\bullet$ Socle sans contagion~: \n{5\,275}~M\euro. \og Plancher défendable.\fg}
\dire{$\bullet$ Ignorance totale sur la direction ($t=1$)~: \n{[6\,858 ; 8\,697]}, largeur
\n{1\,839}, soit \n{34,9\,\%} du socle.}
\dire{$\bullet$ \n{La contagion ajoute entre $+30$ et $+65\,\%$ au socle.}}
\dire{$\bullet$ Le point d'expert (8\,110) est \n{un} point de l'ensemble, \n{pas sa mesure}.}
\dire{$\bullet$ $\rho(W)=\n{0{,}506}$~: la cascade reste sous-critique.}
\dire{$\bullet$ Contrôle croisé~: un second moteur donne 8\,553, soit \n{5\,\%} d'écart.}
\dire{{\color{warm}\textbf{Mots à tenir}~: socle, identification partielle, $t$, rayon
spectral, sous-criticité.}}

\slide{4. À quelle échelle}{le point le plus neuf}
\dire{$\bullet$ \n{\og Ces montants sont un SCR de marché, pas d'entreprise.\fg} À dire tôt
et fort.}
\dire{$\bullet$ Calage secteur mondial~: \n{21,6} incidents TIC matériels par an.}
\dire{$\bullet$ En ne changeant \emph{que} $\lambda$~: 8\,123 $\to$ \n{169}~M\euro.}
\dire{$\bullet$ Les deux lignes du milieu sont des \n{seaux}, définis par un filtre sur la
grandeur même qu'on estime. D'où $\lambda$ estimé comme \n{fonction de la taille}.}
\dire{$\bullet$ \n{Le fait actuariel}~: la moyenne suit $\lambda$ (facteur 270), la
\n{VaR ne suit pas} (facteur 48). \n{Perte unique dominante}, \n{98,8\,\%} des années sans
aucun incident.}
\dire{{\color{warm}\textbf{Mots à tenir}~: élasticité, perte unique dominante.}}

\slide{5. Face au réel}{J5}
\dire{$\bullet$ Au niveau \n{unitaire}~: VaR \n{663}, TVaR \n{1\,752}~M\euro, dans la
fourchette des grandes pertes cyber réelles.}
\dire{$\bullet$ Au niveau \n{annuel}~: de l'ordre du marché cyber mondial. \n{C'est une
échelle de secteur}, ce qui confirme la slide précédente.}
\dire{{\color{warm}\textbf{Mots à tenir}~: VaR, TVaR, et la différence entre les deux.}}

\slide{6. Ramené à une entité}{le renversement}
\dire{$\bullet$ \n{169}~M\euro{} contre \n{450} de forfait, soit \n{0,38} fois.}
\dire{$\bullet$ \n{\og Comparer les niveaux n'a pas de sens\fg}~: le forfait couvre tout
l'opérationnel, moi seulement le cyber.}
\dire{$\bullet$ \n{Ce qui parle, c'est l'insensibilité}~: le forfait vaut 450 pour un profil
exemplaire comme pour un profil défaillant. \n{Il ne bouge pas d'un euro.} Moi j'écarte les
deux de \n{+45\,\%}.}
\dire{$\bullet$ Réserve à donner \n{soi-même}~: le 169 est une \n{borne supérieure} d'ordre
de grandeur, la sévérité est calibrée sur de grandes institutions. \n{L'écart entre états,
lui, est un rapport et ne dépend pas de cette réserve.}}
\dire{{\color{warm}\textbf{Mots à tenir}~: ORSA, Formule Standard, charge opérationnelle.}}

\slide{7. Sur quel point ne pas parier}{J7}
\dire{$\bullet$ \n{91} excès, $\xi=\n{0{,}598}$, IC bootstrap $90\,\%$ \n{[0,32 ; 0,84]}.}
\dire{$\bullet$ L'échelle des postures~: point \n{657}, prédictive \n{645}, robuste
\n{1\,037}, pire cas \n{1\,331}~M\euro.}
\dire{$\bullet$ \n{1.} La prédictive coïncide presque avec le point~: \n{ce n'est pas
l'intégration de l'incertitude qui manquait, c'est la largeur}.}
\dire{$\bullet$ \n{2.} Le niveau $\beta$ est un \n{choix prudentiel explicite, pas un
calcul}. On l'affiche au lieu de le cacher dans un point.}
\dire{{\color{warm}\textbf{Mots à tenir}~: bootstrap, VaR prédictive, VaR robuste, plug-in.}}

\slide{8. Capital libéré}{placer la phrase}
\dire{$\bullet$ Conforme \n{6\,049}, non conforme \n{20\,188}, écart
$\Delta_{\text{DORA}}=\n{14\,139}$~M\euro.}
\dire{$\bullet$ Dont \n{5\,001} d'\n{interaction}, soit \n{35\,\%}~: \n{les canaux ne
s'additionnent pas, la cascade est super-additive}.}
\dire{$\bullet$ Par levier~: fréquence \n{4\,328}, accumulation P4 \n{1\,728}, propagation
$W$ \n{1\,633}, détection \n{1\,448}.}
\dire{$\bullet$ Portage seul (\n{6\,\%}, marge de risque)~: \n{848}~M\euro/an.}
\dire{$\bullet$ \n{1.} Le capital libéré est \n{concentré}~: la fréquence porte \n{47\,\%}
de la somme des leviers. \n{Prioriser, pas étaler.}}
\dire{$\bullet$ \n{2.} Retour sur le \emph{seul} portage~: \n{6 à 35 ans}. \n{Et le portage
est le plancher}, l'essentiel du $\Delta$SCR est de la perte évitée.}
\dire{{\color{warm}\textbf{Mots à tenir}~: super-additivité, marge de risque, Shapley si on
me demande comment j'attribue.}}

\slide{9. Jacobs}{vérification demandée}
\dire{$\bullet$ \n{1.} Ce n'est \n{pas} un lien fréquence-sévérité, c'est une \n{conversion
de sévérité seule}~: $\ln L=7{,}68+0{,}76\ln X+\varepsilon$.}
\dire{$\bullet$ \n{2.} \n{Le capital principal n'en dépend pas}~: il est calculé sur OpRisk,
déjà en dollars. PRC n'est qu'un axe de sensibilité.}
\dire{$\bullet$ \n{3.} La conversion était traitée comme \n{exacte}, ce qui écrasait la
dispersion. En réinjectant le résidu ($\sigma_\varepsilon=\n{0{,}523}$)~: niveau \n{+7\,\%},
queue \n{1,03 $\to$ 1,07}, mais \n{la largeur relative de l'intervalle ne bouge pas}
(1,43 $\to$ 1,41).}
\dire{{\color{warm}\textbf{Mots à tenir}~: résidu, hétéroscédasticité si on pousse.}}

\slide{10. La fréquence}{vérification demandée}
\dire{$\bullet$ $N_j\mid Y\sim\mathcal{P}(\lambda_j m)$, $m=\exp(aY-a^2/2)$. \n{$\lambda$ est
une propriété d'entité, $a$ est un choc commun. Ils ne se lisent pas au même endroit.}}
\dire{$\bullet$ \n{Les mélanger serait une faute}~: la surdispersion d'un panel firme-année
est dominée par l'\n{hétérogénéité de taille}, pas par le choc commun.}
\dire{$\bullet$ Panel (\n{14\,944} obs.)~: Poisson rejeté, $p=\mathbf{1{,}6\times 10^{-30}}$.}
\dire{$\bullet$ Série détrendée~: $a=\n{0,12}$ à $0,30$ contre \n{0,60} posé, \n{au-dessus de
la borne haute} de l'intervalle (0,223).}
\dire{$\bullet$ Erlang~: forme \n{0,63} ajustée, or une Erlang exige une forme \n{entière
$\ge1$}. \n{Le mélange Gamma est le bon objet.}}
\dire{{\color{warm}\textbf{Mots à tenir}~: surdispersion, indice de dispersion, mélange
Gamma, charge systémique.}}

\slide{11. Points bloquants}{clôture}
\dire{$\bullet$ Dépend d'autres~: le \n{registre d'externalisation} (ingestion prête), la
\n{posture} sur le capital à trancher avec Caroline.}
\dire{$\bullet$ Prochains sujets~: l'\n{allocation aux cinq piliers} (Euler pour le niveau,
Shapley pour le surcoût), la \n{relecture complète}.}
\dire{$\bullet$ \n{\og Le niveau bouge dans la bande, l'ordre des piliers ne bouge pas.\fg}
Et elle survit maintenant à un \n{quatrième axe}~: le changement d'échelle.}
\dire{$\bullet$ \n{\og Ne jamais présenter comme mesuré ce qui est posé.\fg} Deux paramètres
viennent de changer de statut, \n{et dans le bon sens}.}

\end{multicols}

\vspace{-2pt}
\begin{alerte}
\footnotesize
\textbf{Les quatre questions à anticiper.}\\[3pt]
\textbf{\og C'est quoi, un placebo ?\fg} $\rightarrow$ \emph{Je casse volontairement la
structure que je prétends détecter, je relance le même calcul, et je vérifie que le signal
disparaît. Ici je permute l'ordre du temps dans chaque entité. S'il survit, ce n'est pas la
structure qui le produisait.} \textbf{Puis enchaîner} : et c'est pour ça que mon $-0{,}33$
n'est pas un échec.\\[3pt]
\textbf{\og Ton $\xi$ à $0{,}595$, pourquoi pas le $1{,}3$ de Hill ?\fg} $\rightarrow$ Hill
suppose une queue de Pareto pure et il est biaisé en échantillon fini. \emph{Et je l'ai
démontré au lieu de l'invoquer : j'ai simulé 2\,000 échantillons sous mon modèle à
$\xi=0{,}595$, appliqué le même Hill, et les six valeurs observées tombent dans l'intervalle
simulé. Le $1{,}3$ est exactement ce qu'un $\xi$ de $0{,}595$ produit.}\\[3pt]
\textbf{\og Pourquoi $g$ n'est pas calibré ?\fg} $\rightarrow$ Il est posé, et la thèse n'en
dépend pas : le capital est croissant en $g$, donc toute correspondance qui respecte
l'\textbf{ordre} produit l'écart, tant que le forfait reste plat. Seule l'amplitude est un
scénario.\\[3pt]
\textbf{\og Et le SCR d'une entité, il vaut combien ?\fg} $\rightarrow$ $169$~M\euro{} en
\textbf{borne supérieure} d'ordre de grandeur, pas en mesure. Ce que j'affirme, c'est l'écart
entre états et l'ordre des piliers.
\end{alerte}

\vspace{2pt}
\begin{alerte}
\footnotesize
\textbf{Trois pièges de vocabulaire.}
\textbf{(i)} Ne jamais dire \og le SCR DORA de telle entité\fg~: il n'existe aucun module
DORA en Formule Standard, c'est un \textbf{besoin de capital ORSA}.
\textbf{(ii)} Ne pas confondre le $\sigma$ d'échelle de la GPD ($57{,}97$~M\euro) avec
l'écart-type d'un test.
\textbf{(iii)} Ne pas confondre $\rho(W)$, qui gouverne le système \emph{simultané}, avec
$R_0$, qui gouverne la cascade d'\emph{incidents}.\\[3pt]
\textbf{Si je ne sais pas} $\rightarrow$ \og Je ne veux pas te répondre de travers, je te le
mets par écrit ce soir.\fg{} \textbf{Puis le faire le jour même.} C'est infiniment mieux
qu'une définition approximative.
\end{alerte}

\end{document}
